\fichatitulo{02 · APLICAÇÃO PARLAMENTAR}{Artigo de conferência · 2024}{KamerRaad}
{\small ROGIERS, Alexander; BUYL, Maarten; KANG, Bo; DE BIE, Tijl.
\textit{KamerRaad: Enhancing Information Retrieval in Belgian National Politics Through Hierarchical Summarization and Conversational Interfaces}.
In: \textit{Machine Learning and Knowledge Discovery in Databases. Research Track and Demo Track}. Cham: Springer, 2024. v. 14948, p. 409--412.
\href{https://doi.org/10.1007/978-3-031-70371-3_30}{DOI: 10.1007/978-3-031-70371-3\_30}.\par}

\campo{Problema e objetivo}
Facilitar a consulta a documentos parlamentares belgas extensos e dispersos, permitindo ao usuário avançar da pergunta às fontes e a explicações contextualizadas.

\campo{Principal contribuição · síntese interpretativa}
Apresenta uma aplicação que combina sumarização hierárquica, metadados e interação conversacional. Resumos de diferentes extensões ajudam a lidar com o limite de contexto, preservando caminhos de acesso aos documentos originais.

\campo{Método}
Descreve coleta, segmentação e sumarização com GPT-4; recuperação com BGE-M3 e geração com GEITje-7B-ultra (seção 3, p. 3--4 do preprint).

\campo{Resultado apresentado}
Apresenta uma aplicação e seu fluxo de consulta às fontes (figuras 1--2). Na versão consultada, não há comparação quantitativa de desempenho; não é possível estimar o ganho da sumarização hierárquica.

\campo{Excertos-chave · seleção literal}
\begin{itemize}
\excerto{Organização documental}{pre-processing step of hierarchical summarization}{p. 2 do preprint, Introdução, parágrafo Contributions.}
\excerto{Rastreabilidade}{summaries and generated responses back to its source documents}{p. 3 do preprint, seção 2.}
\excerto{Retorno dos usuários}{Users can provide explicit, binary feedback}{p. 4 do preprint, seção 3, Implementation Details.}
\end{itemize}

\campo{Limitações declaradas e análise crítica}
\textbf{Autores:} descrevem limites de contexto como desafio e o uso de feedback como melhoria futura (seções 1 e 3).
\textbf{Análise da ficha:} o preprint não apresenta avaliação humana sistemática nem resultados do feedback; isso não prova que avaliações nunca tenham sido realizadas.

\campo{Aproveitamento na dissertação · proposta}
\textbf{Destino:} seção 3 e descrição do artefato na seção 6.
\textbf{Uso:} comparar o tratamento de documentos longos e a ligação entre resposta e fonte. Orienta a pergunta sobre perda de evidência na sumarização, sem justificar sua adoção automática nem comprovar uma lacuna em toda a literatura.

\registro{Conteúdo consultado: \href{https://arxiv.org/pdf/2404.17597v1}{preprint arXiv:2404.17597v1}, p. 1--4, integral. Metadados da publicação conferidos na Springer; excertos usam a paginação do preprint. A bibliografia local diverge em autor e páginas; detalhes no README.}
